\chapter{TEKNİK UYGULAMA}

Bu bölümde Aura Finance uygulamasının temel işlevlerinin teknik gerçekleştirimi ve algoritma detayları sunulmaktadır.

\section{Kullanıcı Kimlik Doğrulama ve Güvenlik}

Uygulama, güvenli bir oturum yönetimi için JSON Web Token (JWT) tabanlı stateless (durum bilgisi saklamayan) bir yapı kullanır.
\begin{itemize}
    \item \textbf{Kayıt Olma:} Kullanıcı şifreleri sunucuya ulaştığında \texttt{bcryptjs} kütüphanesi kullanılarak tuzlanmış (salted) bir şekilde karma (hash) işlemine tabi tutulur.
    \item \textbf{Giriş Yapma:} Başarılı giriş sonrası sunucu, kullanıcının kimlik bilgilerini içeren ve 7 gün geçerliliği olan bir JWT imzalar.
    \item \textbf{Yetkilendirme:} İstemci tarafındaki her hassas istek (harcama ekleme, bütçe görüntüleme vb.), HTTP "Authorization" başlığında bu token'ı taşır.
\end{itemize}

\section{Harcama Yönetimi ve Veri İşleme}

Harcama yönetimi modülü, kullanıcılara esnek bir veri giriş imkanı sunar. Harcamalar PostgreSQL üzerinde \texttt{Expense} tablosunda saklanır. React tarafında \texttt{useExpenses} hook'u, \texttt{React Query} kullanarak verilerin önbelleğe alınmasını ve sunucu ile asenkron senkronizasyonunu yönetir.

\section{Akıllı Veri Giriş Yöntemleri}

Uygulamanın en önemli özelliklerinden biri, veri girişini kolaylaştıran yardımcı teknolojilerdir.

\subsection{OCR ile Makbuz Okuma}
İstemci tarafında çalışan \texttt{Tesseract.js} motoru sayesinde kullanıcılar makbuz fotoğraflarını yükleyebilirler. Uygulama, görüntüden çıkarılan metin içindeki anahtar kelimeleri (örneğin: "Toplam", "Market", "Benzin") analiz ederek harcama miktarını ve kategorisini otomatik olarak tahmin eder.

\subsection{Sesli Komut ile Harcama Ekleme}
Web Speech API kullanılarak gerçekleştirilen bu özellikte, kullanıcının sesli ifadesi (örneğin: "Yemek için 250 lira harcadım") metne dönüştürülür. Sunucu tarafındaki bir düzenli ifade (regex) analizörü, bu metinden sayısal miktarı ve kategori ismini ayıklayarak formu otomatik doldurur.

\section{AI Insights (Yapay Zekâ İçgörüleri)}

Sunucu tarafında çalışan bir analiz motoru, kullanıcının geçmiş harcama verilerini tarayarak aşağıdaki içgörüleri üretir:
\begin{itemize}
    \item \textbf{Abonelik Tespiti:} Belirli aralıklarla tekrarlanan aynı miktardaki harcamaları tespit eder.
    \item \textbf{Anomali Tespiti:} Kullanıcının alışılagelmiş harcama ortalamasının 2.5 katı üzerindeki işlemleri "Olağandışı Harcama" olarak işaretler.
    \item \textbf{Hız Analizi:} Ayın mevcut günündeki harcama hızını kullanarak ay sonu toplam harcama tahmininde bulunur.
\end{itemize}
